%! Describing the Distribution of a Quantitative Variable — Unit 1, Lesson 1.4 — AP Practice
% EFFL, AP Statistics variant: AP-format practice, assigned but NOT scored — the cover's score
% column reads NA for this component. The graded practice is the `homework` component that
% follows it at the back of the packet.
% MC 2 carries the target misconception as distractor (B) ("skewed left because the pile is on
% the left"); MC 3 carries the EK UNC-1.H.7 over-claim. Free-response (a) spirals to Lesson 1.1
% (classify the variable first) and (d) to Lesson 1.3 (check the interval width before you
% believe the shape).
% Page lockstep with ap_practice_key rests on \writelines{n} in the blank matching a terse
% \ansline{} in the key — keep key answers short enough not to wrap past the reserved lines.
% The next-lesson preview is NOT here: it closes the packet, in ../homework/.
%
% THE DISPLAY: 40 commute times, 10-minute intervals from 0: 5 12 10 6 4 2 0 0 1 = 40.
% Right-skewed, with one employee stranded in the 80-90 interval behind a two-interval gap.
\documentclass[10pt]{article}
\usepackage{apstats-article}
\usepackage{apstats-boxes}

\begin{document}

\pageheader{Unit 1, Lesson 1.4}{AP Practice}
% NAMESTRIP: no \namedateperiod here — the name row lives on the cover only.

\vspace{0.15in}

% Authored BYTE-IDENTICALLY in ap_practice/ and ap_practice_key/ — this box is what keeps the
% blank and the key the same number of pages.
\boxguard[8]
\begin{remindbox}
\small
This page is \textbf{not required and not scored} --- your graded homework is the last section
of this packet. These are AP-format reps: work them for the practice, and bring what you get
stuck on to office hours or study hall.
\end{remindbox}

\vspace{0.12in}

\boxguard[24]
\begin{scenariobox}[The Commute Record]{navy}
\small
A company recorded the one-way commute time, in minutes, of each of its 40 employees. The
histogram below groups those 40 commute times into intervals 10 minutes wide.

\vspace{6pt}
\begin{center}
\begin{tikzpicture}[scale=0.95]
  \foreach \i/\c in {0/5, 1/12, 2/10, 3/6, 4/4, 5/2, 6/0, 7/0, 8/1} {
    \fill[navy] ({\i*0.6},0) rectangle ({(\i+1)*0.6},{\c*0.18});
    \draw[white, very thin] ({\i*0.6},0) rectangle ({(\i+1)*0.6},{\c*0.18});}
  \draw[->, thick] (0,0) -- (0,2.7);
  \draw[->, thick] (0,0) -- (5.9,0);
  \foreach \y/\lab in {0/0, 0.9/5, 1.8/10} {
    \draw (-0.07,\y) -- (0.07,\y); \node[left, font=\scriptsize, xshift=-1pt] at (-0.07,\y) {\lab};}
  \foreach \x/\lab in {0/0, 1.2/20, 2.4/40, 3.6/60, 4.8/80} {
    \draw (\x,-0.07) -- (\x,0.07); \node[below, font=\scriptsize] at (\x,-0.1) {\lab};}
  \node[font=\scriptsize] at (2.7,-0.62) {One-way commute time (minutes)};
  \node[font=\scriptsize, rotate=90] at (-0.62,1.2) {Employees};
\end{tikzpicture}
\end{center}
\end{scenariobox}

\vspace{0.14in}

\boxguard[16]
\begin{headlinebox}{goldbg}
\small
\textbf{Section I --- Multiple Choice.} Circle the single best answer. Every answer choice is
about the context, not just the arithmetic --- read them all before choosing.
\end{headlinebox}

\vspace{0.10in}

\begin{enumerate}[leftmargin=1.9em, itemsep=0.13in]

  \item How many of these 40 employees have a one-way commute of 30 minutes or more?
        \begin{enumerate}[label=\textbf{(\Alph*)}, leftmargin=2.2em, itemsep=2pt]
          \item 6
          \item 10
          \item 13
          \item 17
          \item 27
        \end{enumerate}

  \item Which statement best describes the shape of this distribution?
        \begin{enumerate}[label=\textbf{(\Alph*)}, leftmargin=2.2em, itemsep=2pt]
          \item Approximately symmetric and unimodal.
          \item Skewed to the left, because most of the commutes lie on the left side of the
                graph.
          \item Skewed to the right, because the tail of longer commutes stretches out to the
                right.
          \item Approximately uniform, because every interval contains at least a few employees.
          \item Bimodal, because two bars are taller than the rest.
        \end{enumerate}

  \item Which conclusion is \textbf{supported} by this display?
        \begin{enumerate}[label=\textbf{(\Alph*)}, leftmargin=2.2em, itemsep=2pt]
          \item Employees at this company commute longer than employees at other companies.
          \item Commute times for all workers in this city are skewed to the right.
          \item Among these 40 employees, most commute less than 30 minutes, and one commutes
                more than 80 minutes.
          \item The employee with the longest commute should be removed before the distribution
                is described.
          \item Half of these employees commute more than 40 minutes.
        \end{enumerate}

  \item Before concluding that this distribution is skewed to the right, a careful reader should
        first check
        \begin{enumerate}[label=\textbf{(\Alph*)}, leftmargin=2.2em, itemsep=2pt]
          \item whether commute time is a categorical variable.
          \item the interval width used, because a different width can change the appearance of
                the histogram.
          \item whether the bars have gaps between them, as in a bar graph.
          \item whether any employee was counted twice.
          \item nothing --- a histogram reports the shape of the data as a fact.
        \end{enumerate}

\end{enumerate}

\vspace{0.14in}

\boxguard[16]
\begin{headlinebox}{goldbg}
\small
\textbf{Section II --- Free Response.} Show your reasoning. Every answer must be written
\emph{in the context of the problem} --- that is what earns credit on the AP exam.
\end{headlinebox}

\vspace{0.10in}

\begin{enumerate}[leftmargin=1.9em, itemsep=0.16in, resume]

  \item Use the commute record above.

        \vspace{6pt}
        \textbf{(a)} Classify commute time as categorical or quantitative, and then as discrete
        or continuous. Justify from the variable itself, not from the display.\\[4pt]
        \writelines{2}

        \vspace{6pt}
        \textbf{(b)} Describe the shape of this distribution and identify any unusual features,
        in context.\\[4pt]
        \writelines{3}

        \vspace{6pt}
        \textbf{(c)} A manager writes: \emph{``The average commute here is about 25 minutes, so
        our employees are fine.''} Explain what that single number leaves out, and state what
        these data can and cannot tell you about employees at other companies.\\[4pt]
        \writelines{4}

        \vspace{6pt}
        \textbf{(d)} The same 40 commute times are redrawn using intervals 30 minutes wide.
        Explain what happens to the employee who commutes more than 80 minutes, and why the
        interval width has to be checked before the shape is believed.\\[4pt]
        \writelines{3}

\end{enumerate}

\vspace{0.12in}

\boxguard[12]
\begin{extensionbox}
\small
Find a graph in a news article that reports one number as ``the average'' or ``the typical''
value. Say which of the four parts of a description that number leaves out, and what a reader
could be misled into believing because of it.\\[4pt]
\writelines{2}
\end{extensionbox}

\vspace{0.12in}


\end{document}
